% concept01.tex
\section{Variables and Algebraic Expressions \kr{(문자와 식)}}

Welcome to the world of algebra! In this chapter, we'll learn how to use letters like \textbf{x} and \textbf{y} to represent numbers. These letters are called \textbf{variables}, and they help us write expressions and solve equations. Algebra might seem like a new language, but once you understand the rules, it can actually make math easier. Let's explore how we can describe patterns, solve problems, and think logically using simple algebraic tools.

\blockquote{A variable is a symbol, like \(x\) or \(y\), that stands for a number.}

\subsection{What You'll Learn}
\begin{itemize}
    \item Understand what variables are and why we use them
    \item Learn how to write and read algebraic expressions
\end{itemize}

\subsection{Key Ideas}
\begin{itemize}
    \item A \textbf{variable} is a symbol, like $x$ or $y$, that stands for a number.
    \item An \textbf{algebraic expression} is a mix of numbers, variables, and math operations such as $+$, $-$, $\times$, or $\div$.
\end{itemize}

\textbf{Example:}
\[
3x + 5 \quad \text{means “3 times some number, then add 5.”}
\]
If $x = 2$, then:
\[
3x + 5 = 3(2) + 5 = 6 + 5 = 11
\]

\subsection{Everyday Example}
Imagine you earn \$10 per hour. If $h$ is the number of hours you work, then your earnings are written as:
\[
10h
\]
If you work 5 hours:
\[
10(5) = 50
\]

\subsection{Did You Know?}
The word \textbf{algebra} comes from the Arabic term \textit{al-jabr}, which means “reunion of broken parts.” It appeared in a famous math book written over 1,200 years ago by a scholar named Al-Khwarizmi. Many English math words like \textit{algorithm} also come from his name!

\subsection{Tips for Students}
\begin{itemize}
    \item Think of a variable as a mystery box.
    \item Try different values to see how expressions change.
    \item Use colors, arrows, or circles to match terms and operations.
\end{itemize}


% concept02.tex
\section{Simplifying Expressions \kr{(식의 간단화)}}

\blockquote{Like terms have the same variable and exponent.}

\subsection{What You'll Learn}
\begin{itemize}
    \item Combine like terms
    \item Use the distributive property to simplify expressions
\end{itemize}

\subsection{Key Ideas}
\begin{itemize}
    \item \textbf{Like terms} have the same variable and exponent.
    \item The \textbf{distributive property}:
    \[
    a(b + c) = ab + ac
    \]
\end{itemize}

\textbf{Examples:}
\[
2x + 3x = 5x
\]
\[
3(x + 4) = 3x + 12
\]

\subsection{Why It Matters}
Simplifying expressions helps make your work cleaner and your thinking clearer.

\subsection{Did You Know?}
Even NASA engineers and video game developers simplify expressions when building models and writing code.

\subsection{Tips for Students}
\begin{itemize}
    \item Circle or highlight like terms before combining.
    \item Be careful with minus signs.
    \item Go one step at a time.
\end{itemize}


% concept03.tex
\section{Solving One-Step Equations \kr{(1단계 방정식)}}

\blockquote{Use inverse operations to isolate the variable. Do the same to both sides.}

\subsection{What You'll Learn}
\begin{itemize}
    \item Solve equations in one step
    \item Understand inverse operations
\end{itemize}

\subsection{Key Ideas}
\begin{itemize}
    \item Use \textbf{inverse operations} to isolate the variable.
    \item Do the same thing to both sides of the equation.
\end{itemize}

\textbf{Examples:}
\[
x + 4 = 9 \Rightarrow x = 5
\]
\[
5x = 20 \Rightarrow x = 4
\]

\subsection{Real-World Example}
If \$20 is 5 times the cost of one notebook:
\[
5x = 20 \Rightarrow x = 4
\]

\subsection{Did You Know?}
Solving equations is like solving a riddle or puzzle.

\subsection{Tips for Students}
\begin{itemize}
    \item Ask: “What's happening to the variable?” Then do the opposite.
    \item Keep both sides equal.
    \item Plug your answer back in to check.
\end{itemize}


% concept04.tex
\section{Solving Multi-Step Equations \kr{(여러 단계 방정식)}}

\blockquote{First, simplify both sides. Then, isolate the variable step by step.}

\subsection{What You'll Learn}
\begin{itemize}
    \item Solve equations with more than one operation
    \item Use combining like terms and distribution
\end{itemize}

\subsection{Key Ideas}
\begin{itemize}
    \item First, simplify both sides.
    \item Then, isolate the variable step by step.
\end{itemize}

\textbf{Example:}
\begin{align*}
2(x + 3) &= 14 \\
2x + 6 &= 14 \\
2x &= 8 \\
x &= 4
\end{align*}

\subsection{How This Helps}
Solving multi-step equations builds confidence and logical thinking.

\subsection{Did You Know?}
Multi-step equations are used in budgeting, engineering, and cooking.

\subsection{Tips for Students}
\begin{itemize}
    \item Don't rush—write each step clearly.
    \item Simplify both sides before solving.
    \item Check your solution.
\end{itemize}


% concept05.tex
\section{Inequalities \kr{(부등식)}}

\blockquote{If you multiply or divide by a negative number, flip the inequality sign.}

\subsection{What You'll Learn}
\begin{itemize}
    \item Solve basic inequalities
    \item Graph answers on a number line
\end{itemize}

\subsection{Key Ideas}
\begin{itemize}
    \item Inequality signs:
    \begin{itemize}
        \item $>$ greater than
        \item $<$ less than
        \item $\geq$ greater than or equal to
        \item $\leq$ less than or equal to
    \end{itemize}
    \item If you multiply or divide by a negative number, \textbf{flip the inequality sign}.
\end{itemize}

\textbf{Examples:}
\[
x - 3 > 2 \Rightarrow x > 5
\]
\[
-2x < 6 \Rightarrow x > -3
\]

\subsection{Everyday Example}
A sign says: “You must be under 18 to enter.” That means:
\[
x < 18
\]

\subsection{Did You Know?}
Inequalities are used in business, programming, and transportation.

\subsection{Tips for Students}
\begin{itemize}
    \item Treat inequalities like equations—unless multiplying/dividing by a negative.
    \item Graph solutions:
    \begin{itemize}
        \item Open circles ($\circ$) for $<$ or $>$
        \item Closed circles ($\bullet$) for $\leq$ or $\geq$
    \end{itemize}
    \item Test a number to check your solution.
\end{itemize}